\section{Numerical Simulation}

%================================================
\begin{frame}{Finite Element Solution of the PDE}

	\begin{block}{Weak Form}
		\[
			\int_\Omega
			\frac{\partial c}{\partial t}v\,dx
			+
			D
			\int_\Omega
			\nabla c \cdot \nabla v\,dx
			=
			\int_\Omega
			\alpha G_{R_t}(\mathbf{x}-\mathbf X_t)v\,dx
		\]
	\end{block}

	\begin{alertblock}{Boundary Condition}
		Reflective (no-flux) boundary:
		\[
			\frac{\partial c}{\partial n}=0
			\qquad \text{on } \partial\Omega
		\]
	\end{alertblock}

	\begin{itemize}
		\item Finite-element discretization.
		\item Implemented using FEniCS.
	\end{itemize}

\end{frame}

%================================================
\begin{frame}{Numerical Integration of the SDE}
	\begin{block}{Drift term}
		\[
			\mathbf{A}(\mathbf{X}_t,t)
			=
			-\frac{\beta R_t^3}{\gamma_t\delta^3}
			\int_{\Omega}
			G_{R_t}(\mathbf{x}-\mathbf{X}_t)\,
			\nabla c_t(\mathbf{x})\,d\mathbf{x}
		\]
	\end{block}

	\begin{block}{Euler--Maruyama Scheme}

		\[
			\mathbf X^{n+1}
			=
			\mathbf X^{n}
			+
			\mathbf A(\mathbf X^n,t_n)\Delta t
			+
			\frac{\sqrt{\sigma}}{\eta_t}
			\sqrt{\Delta t}\,
			\eta^n
		\]

		\[
			\eta^n \sim \mathcal N(0,I)
		\]

	\end{block}

	\begin{itemize}
		\item Deterministic drift from chemical gradients.
		\item Random fluctuations modeled by Gaussian noise.
		\item Generates stochastic droplet trajectories.
	\end{itemize}

\end{frame}
%================================================
\begin{frame}{Simulation Workflow}

	\begin{columns}[T]

		\column{0.45\textwidth}

		\centering
		\begin{tikzpicture}[node distance=0.8cm,
				every node/.style={
						draw,
						rounded corners,
						align=center,
						minimum width=3.5cm,
						minimum height=0.6cm}]

			\node (init) {Initialize $c(\mathbf{x},0)$ and $\mathbf{X}_0$};
			\node (pde) [below=of init] {Solve PDE for $c(\mathbf{x},t)$};
			\node (grad) [below=of pde] {Compute $\nabla c$};
			\node (sde) [below=of grad] {Update $\mathbf{X}_t$ (Euler--Maruyama)};
			\node (src) [below=of sde] {Update source term};

			\draw[->] (init)--(pde);
			\draw[->] (pde)--(grad);
			\draw[->] (grad)--(sde);
			\draw[->] (sde)--(src);
			\draw[->,bend left=40] (src.east) to node[right] {\small repeat} (pde.east);

		\end{tikzpicture}

		\column{0.55\textwidth}

		\begin{block}{Key Steps}
			\begin{itemize}
				\item Evolve the chemical field using the diffusion PDE.
				\item Compute the gradient-induced drift.
				\item Update droplet position via Euler--Maruyama and refresh the source term.
			\end{itemize}
		\end{block}

		\begin{alertblock}{Outcome}
			Simultaneous evolution of the chemical field and droplet trajectory generates the simulated paths.
		\end{alertblock}

	\end{columns}

\end{frame}
% %================================================
% \begin{frame}{Simulation Algorithm}

% 	\begin{enumerate}
% 		\item Initialize the droplet position \(\mathbf{X}_0\) and concentration field \(c(\mathbf{x},0)\).
% 		\item Solve the diffusion equation for the concentration field using finite element discretization.
% 		\item Compute the chemical-gradient-induced drift term from the concentration field.
% 		\item Update the droplet position using the Euler--Maruyama scheme.
% 		\item Update the source term according to the new droplet position.
% 		\item Repeat the procedure for all time steps.
% 	\end{enumerate}

% 	\vspace{0.4cm}

% 	\begin{block}{Outcome}
% 		This coupled procedure evolves both the concentration field and the droplet trajectory simultaneously.
% 	\end{block}

% \end{frame}

%================================================
\begin{frame}{Comparison Between Experiment and Simulation}

	\begin{alertblock}{Ensemble Analysis Procedure}

		\begin{itemize}

			\item MSD, VACF, and OCF were computed for each retained trajectory.

			\item Lag times were restricted to approximately
			      $40\%$ of the trajectory length to ensure reliable statistics.

			\item Ensemble observables were obtained by averaging over all retained droplets.

			\item MSD exponents were obtained from
			      $\mathrm{MSD}(\tau)\sim\tau^\alpha$
			      using the first $30\%$ of the available lag-time range.
			\item Exponential fits to VACF and OCF were used
			      to estimate characteristic persistence times
			      $\tau_p$ and $\tau_r$.
			\item Exponential fits
			      $C_v(\tau)\sim e^{-\tau/\tau_p}$ and
			      $C_\theta(\tau)\sim e^{-\tau/\tau_r}$
			      yielded $\tau_p$ and $\tau_r$.

		\end{itemize}

	\end{alertblock}


\end{frame}
%================================================
